%%\documentclass{article}
%%\documentclass[12pt,a4paper]{article}
\documentclass[12pt,reqno,a4paper]{amsart}

% --- preamble --- 


% 1. Encoding and Fonts (standard)
\usepackage[utf8]{inputenc}
\usepackage[T1]{fontenc}

%%\usepackage{amsaddr}
% 2. Math and standard AMS packages
\usepackage{amsmath, amsfonts, amssymb, amsthm}
%%\usepackage[hyper]{amsbib}

% more packages

\usepackage{listings}
\lstset{ %
	breaklines=true, %
	backgroundcolor=\color{white}, %
	commentstyle=\color{green}, %
	keywordstyle=\color{blue}\textbf,  %
stringstyle=\color{orange}, %
identifierstyle=\color{red}\itshape,  %
}
\usepackage{color}
\usepackage{imakeidx}
% generate an index file:
\makeindex


% 3. Metadata and Linking (Load these LAST)
\usepackage[hidelinks]{hyperref} % 'hidelinks' removes the ugly red/green boxes
\usepackage{hyperxmp}

\hypersetup{
%%    pdftitle={Paper Title
    pdfauthor={Author Name},
    pdfcopyright={Licensed under CC BY 4.0},
    pdflicenseurl={https://creativecommons.org/licenses/by/4.0/}
}

% 4. Optional: Simple Footer License
\usepackage{fancyhdr}
\pagestyle{fancy}
\fancyhf{}
\renewcommand{\headrulewidth}{0pt} % Remove the horizontal line at the top
\fancyfoot[C]{\thepage}
\fancyfoot[L]{\footnotesize CC BY 4.0}

% --- article metadata ( titlepage info) ---

\title[paper title]{Paper Title}
%\author[Lastname1]{Firstname1 Lastname1}
% Source - https://tex.stackexchange.com/a/311688
% Posted by Dai Bowen, modified by community. See post 'Timeline' for change history
% Retrieved 2026-01-21, License - CC BY-SA 3.0

\author[F1. LastName1]{FirstName1 LastName1$^1$}
\address{Address1}
\thanks{$^1$Address1, Email: \href{mailto:first1.last1@domain1}{first1.last1@domain1}}
\author[F2. LastName2]{FirstName2 LastName2$^2$}
\address{Address2}
\thanks{$^2$Address2, Email: \href{mailto:first2.last2@domain2}{first2.last2@domain2}}

%  Standard amsart: email and address print at the very end
% \email{\href{mailto:mame@mail.domain}{mailto:name@mail.domain}}
% Optional: Use \thanks if you want the email on page 1
% %\thanks{Email: email@example.com}
% \thanks{Email: \href{mailto:mame@mail.domain}{mailto:name@mail.domain}}
%\thanks{Email: \href{mame@mail.domain}{mailto:name@mail.domain}}
%\address{The address}
%\date{\today}


% --- document body ---
\begin{document}

\begin{abstract}

Write the paper abstract here

\end{abstract}
%% note: for amsart, abstact must precede \maketitle
\maketitle

\pagestyle{headings}


\cleardoublepage
\phantomsection % This creates a precise anchor for the PDF link
\addcontentsline{toc}{section}{Contents}
\tableofcontents

\addcontentsline{}{section}{Introduction}
\section*{Introduction}


This aper is  written in {\LaTeX}, (see \cite{latex2e} and see  {\TeX} \cite{texbook} for {\TeX}.)

It uses the \texttt{amsart} document class.





\section{a first section}
\label{sec:first}

This  is the first section

\subsection{a subsection}

Paragraphs are automatically created from blocks of text separated by one or more (tradittionally two or more)
blank lines.


Some more text goes here in a new paragraph.




\subsection{another subsection}


\subsection{ath formulae}

Writing math\index{math} formulae\index{formulae} is rather easy. 

Here is an example with a numbered, displayed equation:


\begin{equation}
\label{eq:example}
\sum_{i=1}^{n}i = \frac{n(n+1)}{2}
\end{equation}


 Alternatively, formulae can be written inline (without numbering). For example $a^2 + b^2 = c^2$.

\section{second section}

\subsection{a subsection}


\subsection{another subsection}


\section{third section}

%%\begin{verbatim}

\begin{lstlisting}
 # pdflatex doc
 # pdflatex doc
 # bibtex doc
 # makeindex doc
 # pdflatex doc
\end{lstlisting}
%%\end {verbatim}


\section{fourth section}







\section{final section}

closing words.



%%\section*{Index}
% load and print  the index:
\printindex
\addcontentsline{}{section}{Index}



\bibliography{doc}{}
\bibliographystyle{amsplain}

\addcontentsline{}{section}{References}


\end{document}
